% ============================================================
%  ch01b  §1.2 驯顺分歧的局部几何朗兰兹对应
%  原文对应：ICM-arXiv.tex 658--748 行
%  编号：公式 (1.7)(1.8)(1.9)；注 1.2.1；定理 1.2.2；注 1.2.3；
%        定理 1.2.4；注 1.2.5；注 1.2.6
% ============================================================

\subsection{驯顺分歧的局部几何朗兰兹对应}\label{Tamely local}

先回顾经典理论。设 $F$ 是 $\bQ_p$ 的有限扩张，或同构于 $\bF_q(\!(\varpi)\!)$；为简单起见，设 $G$ 可延拓为 $\mO$ 上的连通约化群。（事实上，本小节的结果在 $F$ 的驯顺分歧扩域上分裂的拟分裂群情形下，都以适当形式成立。）此外，固定 $G$ 在 $\mO$ 上的一个钉 $(B,T,e)$。

经典局部朗兰兹纲领的目标，是用 Galois 表示来分类 $G(F)$ 的（光滑）不可约表示（系数取 $\bC$）。从这个观点看，Satake 同构 \eqref{Sat isom} 给出了 $G(F)$ 不可约非分歧表示的分类，即那些含有 $G(\mO)$-不动非零向量的表示：这类表示与 $H^{\cl}_{G(\mO)}\otimes\bC$ 的单模一一对应，而由 Satake 同构 \eqref{Sat isom}，后者又由 $\CG$ 中的半单 $\hat{G}$-共轭类参数化。（对不可约非分歧表示 $\pi$，与之对应的 $\CG$ 中的 $\hat{G}$-共轭类通常称为 $\pi$ 的 Satake 参数。）

下一类重要的不可约表示，是那些含有 Iwahori 子群不动非零向量的表示。例如，在模 $\varpi$ 约化映射 $G(\mO)\to G(k)$ 下，$B(k)\subset G(k)$ 的原像 $I$ 就是 $G(F)$ 的一个 Iwahori 子群。与非分歧情形一样，$I$-双不变的 $\bZ$ 值函数构成一个 $\bZ$-代数 $H^{\cl}_{I}$，其乘法由卷积 \eqref{E:conv prod alge} 给出（Haar 测度 $dg$ 取使 $I$ 的体积为 $1$ 者）；而 $G(F)$ 的含有非零 $I$-不动向量的不可约表示，与 $(H^{\cl}_I\otimes\bC)$-单模一一对应。

与 Satake 同构类似，Kazhdan--Lusztig 用与 $\hat{G}$ 相伴的几何对象描述了 $H^{\cl}_{I}\otimes\bC$。

以 $\hat{U}\subset\hat{B}$ 记 $\hat{B}$ 的幂幺根。自然态射 $\hat{U}/\hat{B}\to \hat{G}/\hat{G}$ 通常称为 Springer 消解。令
\[
S^{\un}_{\hat{G}}=(\hat{U}/\hat{B})\times_{\hat{G}/\hat{G}}(\hat{U}/\hat{B}),
\]
称之为 $\hat{G}$ 的（幂幺）Steinberg 叠\footnote{由于 $\hat{U}/\hat{B}\to \hat{G}/\hat{G}$ 不是平坦的，纤维积须在导出意义下取，故 $S^{\un}_{\hat{G}}$ 应理解为导出代数叠。}。在 $\bC$ 上，借助指数映射把 $\hat{U}_\bC$ 与其 Lie 代数等同，可以得到 $\hat{U}_\bC$ 上的 $\bG_{m,\bC}$-作用，从而得到 $S^{\un}_{\hat{G},\bC}$ 上的 $\bG_{m,\bC}$-作用。于是可以取商叠 $S^{\un}_{\hat{G},\bC}/\bG_{m,\bC}$。下文对 $\bC$ 上的（有限型的）Artin 叠 $X$，以 $K(X)$ 记 $X$ 上凝聚层（$\infty$-）范畴的 $K$-群。

Kazhdan--Lusztig \cite{Kazhdan-Lusztig-Iwahori} 在 $G$ 分裂且中心连通的假设下，构造了典范同构（需先取定 $q$ 的一个平方根 $\sqrt{q}$）
\begin{equation}\label{KL isomorphism}
K(S^{\un}_{\hat{G},\bC}/\bG_{m,\bC})\otimes_{K(\bB\bG_{m,\bC})}\bC\cong H^{\cl}_I\otimes\bC,
\end{equation}
其中映射 $K(\bB\bG_{m,\bC})\to \bC$ 把与 $\bG_{m,\bC}$ 的重言表示相对应的类送到 $\sqrt{q}$。此外，该同构诱导出 Bernstein 同构
\begin{equation}\label{Bernstein isomorphism}
K(\bB \hat{G}_\bC)\otimes\bC\cong Z(H^{\cl}_I\otimes\bC),
\end{equation}
其中 $Z(H^{\cl}_I\otimes\bC)$ 是 $H^{\cl}_I\otimes\bC$ 的中心，而映射 $K(\bB \hat{G}_\bC)\to K(S^{\un}_{\hat{G},\bC}/\bG_{m,\bC})$ 由自然投影 $S^{\un}_{\hat{G}}/\bG_m\to \bB\hat{G}$ 诱导。

\begin{rmk}
用涉及 $\hat{G}$ 的几何来刻画 $\bZ$-代数 $H^{\cl}_I$，将是很有意思的问题；这将推广 \cite{Zhu-Integral-Satake} 中的整系数 Satake 同构。
\end{rmk}

Kazhdan--Lusztig 同构 \eqref{KL isomorphism} 也容许范畴化，通常称为 Bezrukavnikov 等价：它给出仿射 Hecke 范畴的两种实现。与前面一样，转入几何理论时允许 $F$ 稍更一般，如 \S\ref{S: geometric Satake} 中那样；同时设 $G$ 可延拓为 $\mO$ 上的连通约化群，并固定 $G$ 在 $\mO$ 上的一个钉。

设 $L^+G\to G_k$ 为自然的“模 $\varpi$ 约化”映射，以 $\Iw\subset L^+G$ 记 $B_k\subset G_k$ 的原像。这是 $I$ 的几何类比。于是如 \S\ref{S: geometric Satake} 中讨论的非分歧情形，可以定义 Iwahori 局部 Hecke 叠 $\Hk_{\Iw}= \Iw\backslash LG/\Iw$ 与幺半范畴
$\Shv_{c}(\Hk_{\Iw}\otimes\overline k,\La)\subset\Shv(\Hk_{\Iw}\otimes\overline k,\La)$。范畴 $\Shv_{c}(\Hk_{\Iw}\otimes\bar k,\La)$ 可看作 $H^{\cl}_I$ 的范畴类比，通常称为仿射 Hecke 范畴。

回忆 $\breve F=W_\mO(\overline k)[1/\varpi]$。$F$ 的惯性群 $I_F:=\Ga_{\breve F}$ 有驯顺商 $I_F^t$，同构于 $\prod_{\ell\neq p} \bZ_\ell(1)$。

\begin{thm} \label{T: Bezrukavnikov equivalence}
对驯顺惯性群 $I_F^t$ 的任一拓扑生成元 $\tau$ 的选取，存在幺半 $\infty$-范畴之间的典范等价
\[
\on{Bez}^{\un}_G:   \Coh(S^{\un}_{\hat{G},\La})\cong \Shv_{c}(\Hk_{\Iw}\otimes\overline k,\La).
\]
\end{thm}

事实上，Bezrukavnikov 在 $F=k(\!(\varpi)\!)$ 时于 \cite{Bezrukavnikov-Affine-Hecke} 中证明了这一等价；Yun 与作者由等特征情形推出混合特征情形。一个有意思的问题是：\cite{Scholze-Berkeley, Fargues-Scholze-Geometrization} 引入的新技术能否直接给出该等价在混合特征情形的证明。（这方面的若干进展见 \cite{AGLR}。）

\begin{rmk}
再次指出，对算术应用而言，需要用对偶群一侧来描述 $\Ga_k$ 在 $\Shv_{c}(\Hk_{\Iw}\otimes\bar k,\La)$ 上的作用。讨论见 \cite{Bezrukavnikov-Affine-Hecke, Hemo-Zhu-Unipotent}。
\end{rmk}

我们说明定理 \ref{T: Bezrukavnikov equivalence} 的证明（在 $F=k(\!(\varpi)\!)$ 时）中的一个重要环节。存在 $\mO$ 上光滑仿射群概形 $\mG$（称为 Iwahori 群概形），使得 $\mG\otimes F=G$ 且 $L^+\mG=\Iw$。于是有 $D$ 上的局部 Hecke 叠 $\Hk_{\mG,D}$，它类似于 \S\ref{S: geometric Satake} 末尾讨论的 $\Hk_{G,D^S}$（此处 $S=\{1\}$）。此外 $\Hk_{\mG,D}|_{D^*}\cong \Hk_{G,D}|_{D^*}$，且 $\Hk_{\mG,D}|_0=\Hk_{\Iw}$，其中 $0\in D$ 是闭点。

取邻近圈便得到
\begin{equation}\label{E: central sheaf}
\mZ: \Coh(\bB\hat{G}_\La)^\heart\xrightarrow{\Sat_{\{1\}}} \Perv(\Hk_{\mG,D}|_{D^*_{\bar k}},\La)\xrightarrow{\Psi} \Perv(\Hk_{\Iw}\otimes \bar k,\La).
\end{equation}
这就是 Gaitsgory 的中心函子 \cite{Gaitsgory-Central-Sheaf, Zhu-Coherence-Conjecture}，可视为 \eqref{Bernstein isomorphism} 的范畴化。这一构造的动机来自研究志村簇 mod $p$ 几何时提出的 Kottwitz 猜想，相关讨论见 \S\ref{S: local model}。

定理 \ref{T: Bezrukavnikov equivalence} 可以推广到驯顺层。考虑交换图
\[
\hat{G}/\hat{G}\leftarrow \hat{B}/\hat{B}\xrightarrow{q_{\hat{B}}} \hat{T}/\hat{T},
\]
其中左态射是通常的 Grothendieck--Springer 消解。设 $\chi$ 是 $\hat{T}/\hat{T}$ 的一个 $\La$-点，其中 $\La$ 是 $\bQ_\ell$ 的有限扩张。令 $(\hat{B}/\hat{B})_\chi=q_{\hat{B}}^{-1}(\chi)$，并令
\[
S^{\chi}_{\hat{G},\La}:= (\hat{B}/\hat{B})_\chi\times_{\hat{G}/\hat{G}} (\hat{B}/\hat{B})_\chi.
\]
注意当 $\chi=1$ 时它化归为 $S^{\un}_{\hat{G},\La}$。

另一方面，（挠的）$\La$-点 $\chi\in \hat{T}/\hat{T}$ 定义了 $\Iw\otimes\overline{k}$ 上的一维特征层 $\mL_\chi$。于是可以定义 $LG_{\bar k}$ 上双 $(\Iw,\mL_\chi)$-等变可构造层的幺半范畴，记作 $\Shv_{\cons}({}_\chi(\Hk_{\Iw})_{\chi},\La)$。当 $\chi=1$ 时 $\mL_\chi$ 就是 $\Iw$ 上的平凡特征层，此即仿射 Hecke 范畴 $\Shv_c(\Hk_{\Iw}\otimes\overline k,\La)$。定理 \ref{T: Bezrukavnikov equivalence} 的如下推广由 \cite{Bezrukavnikov-Affine-Hecke} 提出，它将在与 Dhillon--Li--Yun 的合作文章 \cite{DLYZ-Endoscopy} 中证明。

\begin{thm} \label{T: Monodromic Bezrukavnikov equivalence}
设 $\cha F=\cha k$。存在典范的幺半等价
\[
\on{Bez}^{\chi}_{\hat{G}}: \Coh(S^{\chi}_{\hat{G},\La})\cong \Shv_{c}({}_\chi(\Hk_{\Iw})_{\chi},\La).
\]
\end{thm}

\begin{rmk}\label{R: int Bez}
为定理 \ref{T: Bezrukavnikov equivalence} 与定理 \ref{T: Monodromic Bezrukavnikov equivalence} 中的等价建立 $\bZ_\ell$-层版本，是重要的。
\end{rmk}

\begin{rmk}\label{2-categorical LGL}
驯顺分歧之外的局部几何朗兰兹对应尚未被完全理解，尽管某些野分歧已在具体问题中出现（例如 \cite{HNY-Kloosterman, Zhu-FG-vs-HNY}）。人们普遍相信，一般的局部几何朗兰兹应以 2-范畴式的陈述来表述：它预测，在 $LG$ 上（适当定义的）层范畴作用下，模范畴构成的 2-范畴，与局部几何朗兰兹参数叠上的范畴构成的 2-范畴等价。确切表述超出本文范围；粗略地说，它蕴含（并且多少等价于说）$LG$ 在适当选取的“水平”下的 Hecke 范畴（Morita）等价于某个形如 $X\times_{Y}X$ 的叠上凝聚层范畴，其中 $Y$ 与局部几何朗兰兹参数的模空间密切相关。
\end{rmk}
